\setcounterpageref{ex}{0}
\TaoPhieuKiemTra{02}

\setcounter{ex}{0}
\setcounter{bt}{0}

\textbf{\large I. PHẦN TRẮC NGHIỆM (3,0 điểm)}

{\itshape Khoanh tròn vào chữ cái đặt trước câu trả lời đúng. Mỗi câu đúng được 0,5 điểm.}

\begin{ex} % Câu 1
Số gồm $7$ chục và $4$ đơn vị viết là
\choice
{$47$}
{$704$}
{\True $74$}
{$70$}
\loigiai{
Số gồm $7$ chục và $4$ đơn vị là $74$.\\
\textbf{Đáp án đúng: C}
}
\end{ex}

\begin{ex} % Câu 2
Trong phép tính $68 - 25 = 43$, số $68$ được gọi là
\choice
{\True Số bị trừ}
{Số trừ}
{Hiệu}
{Số hạng}
\loigiai{
Số đứng trước dấu trừ gọi là Số bị trừ.\\
\textbf{Đáp án đúng: A}
}
\end{ex}

\begin{ex} % Câu 3
Đơn vị lít được viết tắt là
\choice
{$\text{kg}$}
{\True $l$}
{$\text{m}$}
{$\text{dm}$}
\loigiai{
Lít được viết tắt là $l$.\\
\textbf{Đáp án đúng: B}
}
\end{ex}

\begin{ex} % Câu 4
Người ta thường dùng chữ cái nào để đặt tên cho một điểm?
\choice
{Chữ cái in thường ($a, b, c, \dots$)}
{Chữ số ($1, 2, 3$)}
{Kí hiệu hình vẽ}
{\True Chữ cái in hoa ($A, B, C, \dots$)}
\loigiai{
Theo quy ước toán học, tên điểm dùng chữ in hoa.\\
\textbf{Đáp án đúng: D}
}
\end{ex}

\begin{ex} % Câu 5
Điền dấu thích hợp: $14 - 6 \dots 15 - 7$.
\choice
{\True $=$}
{$>$}
{$<$}
{Không so sánh được}
\loigiai{
$14 - 6 = 8$ và $15 - 7 = 8$. Vì $8 = 8$ nên điền dấu $=$.\\
\textbf{Đáp án đúng: A}
}
\end{ex}

\begin{ex} % Câu 6
Tính: $14 - 8 + 7 = \dots$
\choice
{$12$}
{$14$}
{\True $13$}
{$11$}
\loigiai{
$14 - 8 = 6$; $6 + 7 = 13$.\\
\textbf{Đáp án đúng: C}
}
\end{ex}

\vspace{4mm}
\textbf{\large II. PHẦN TỰ LUẬN (7,0 điểm)}

{\itshape Thí sinh trình bày chi tiết bài làm vào giấy kiểm tra.}

\begin{bt}[1,0 điểm]
Tính nhẩm:
\begin{multicols}{2}
\begin{enumerate}[a)]
  \item $8 + 6$
  \item $7 + 7$
  \item $14 - 7$
  \item $16 - 8$
\end{enumerate}
\end{multicols}
\loigiai{
\begin{itemize}
  \item $8 + 6 = 14$
  \item $7 + 7 = 14$
  \item $14 - 7 = 7$
  \item $16 - 8 = 8$
\end{itemize}
}
\end{bt}

\begin{bt}[1,5 điểm]
Đặt tính rồi tính:
\begin{multicols}{3}
\begin{enumerate}[a)]
  \item $27 + 38$
  \item $79 - 27$
  \item $45 + 52$
\end{enumerate}
\end{multicols}
\loigiai{
\begin{itemize}
  \item $27 + 38 = 65$
  \item $79 - 27 = 52$
  \item $45 + 52 = 97$
\end{itemize}
}
\end{bt}

\begin{bt}[1,0 điểm]
Tính giá trị biểu thức kèm đơn vị đo:
\begin{enumerate}[a)]
  \item $18\text{ kg} + 9\text{ kg} - 11\text{ kg}$
  \item $27\text{ }l - 13\text{ }l + 8\text{ }l$
\end{enumerate}
\loigiai{
\begin{enumerate}[a)]
  \item $18\text{ kg} + 9\text{ kg} - 11\text{ kg} = 27\text{ kg} - 11\text{ kg} = 16\text{ kg}$.
  \item $27\text{ }l - 13\text{ }l + 8\text{ }l = 14\text{ }l + 8\text{ }l = 22\text{ }l$.
\end{enumerate}
}
\end{bt}

\begin{bt}[1,5 điểm]
Bạn Nam có $17$ quyển vở. Bạn Huy có ít hơn bạn Nam $6$ quyển vở. Hỏi bạn Huy có bao nhiêu quyển vở?
\loigiai{
Số quyển vở bạn Huy có là:
\[ 17 - 6 = 11 \text{ (quyển vở)} \]
\textbf{Đáp số:} $11$ quyển vở.
}
\end{bt}

\begin{bt}[1,0 điểm]
Một thùng chứa đầy $37\text{ }\text{kg}$ gạo nếp. Buổi sáng cửa hàng đã bán ra $14\text{ }\text{kg}$. Hỏi trong thùng còn lại bao nhiêu ki-lô-gam gạo nếp?
\loigiai{
Trong thùng còn lại số ki-lô-gam gạo nếp là:
\[ 37 - 14 = 23\text{ }\text{kg} \]
\textbf{Đáp số:} $23\text{ }\text{kg}$ gạo nếp.
}
\end{bt}

\begin{bt}[1,0 điểm]
Điền tiếp hai số thích hợp vào chỗ chấm để hoàn thành quy luật dãy số: $3, 6, 9, 12, \dots, \dots$
\loigiai{
Quy luật: Mỗi số sau bằng số trước cộng thêm $3$.\\
Số tiếp theo: $12 + 3 = 15$.\\
Số tiếp theo nữa: $15 + 3 = 18$.\\
\textbf{Đáp số:} $15$ và $18$.
}
\end{bt}
